% SOURCE_AUTHORITY: sga5_fr_workpass.tex, lines 11057--11100 (LF-normalized unit).
% SOURCE_SHA256: 791F4EFFC5E02832D5D77ED03518C8156D6F07E4C8238B03545DB93D883FBB28
Enunciemos ahora la «fórmula de autointersección» en el anillo de Chow.

\begin{statement}{Teorema 9.2 (Mumford)}
Sean $k$ un cuerpo separablemente cerrado e $i:Y\hookrightarrow X$ una inmersión cerrada, de codimensión pura $d$, de $k$-esquemas cuasiproyectivos y lisos. Entonces, denotando por $N$ el haz normal de $i$, se tiene, para todo $y\in C^*(Y)$, la igualdad
\[
i^*i_*(y)=y\,c_d(\check N)
\]
en $C^*(Y)$.
\end{statement}

Para la demostración necesitaremos cierto número de notaciones y esquemas auxiliares. Embebamos $X$ en $Z=X\times_k\mathbb P^1_k$ mediante el morfismo $u:t\mapsto(t,0)$, de modo que se tiene un diagrama cartesiano
\[
\begin{tikzcd}
\widehat N \arrow[r,"j_1"] \arrow[d,"g_1"'] & Z' \arrow[d,"f_1"]\\
Y \arrow[r,"i_1"'] & Z ,
\end{tikzcd}
\]
donde $Z'$ denota el esquema obtenido por la explosión de $Z$ a lo largo de $Y$, y $\widehat N=\mathbb P_Y(N\oplus\mathcal O_Y)$ es la compleción proyectiva del fibrado vectorial $V(N)$. Como en el enunciado de 9.1, se denota por $s:Y\hookrightarrow\widehat N$ el morfismo obtenido componiendo la sección cero de $V(N)$ con la inclusión canónica $V(N)\subset\widehat N$.

La inclusión $Y\subset X$ permite identificar $Y\times\mathbb P^1$ con un subesquema cerrado de $X\times\mathbb P^1$. Se denota por $W'$ la «transformada estricta de $Y\times\mathbb P^1$ por $f_1$», es decir, el subesquema de $Z'$ obtenido por la explosión de $Y\times\mathbb P^1$ a lo largo de $Y=Y\times0$. Es claro que $W'\cap\widehat N$ se identifica con la sección cero de $\widehat N$ y que, como $Y$ es de codimensión $1$ en $Y\times\mathbb P^1$, el morfismo canónico $f_2:W'\to Y\times\mathbb P^1$ es un isomorfismo. Las demás notaciones relativas a $W'$ quedan fijadas por el diagrama canónico $(D_1)$.
\[
\begin{tikzcd}[column sep=3.2em,row sep=2.4em,cells={nodes={inner sep=1pt}}]
\text{(D}_1\text{)}\arrow[from=1-3,to=5-2,"g_1"']
& & \widehat N \arrow[rr,"j_1"] & & Z' \arrow[ddd,"f_1"] \\
& Y \arrow[ur,"s"] \arrow[rr,"\delta"'] \arrow[ddd,"\mathrm{id}"'] & &
W' \arrow[ur,"\beta"'] \arrow[ddd,"f_2"'] & \\
& & & & \\
& & & & X\times\mathbb P^1=Z \\
& Y \arrow[rr,"\gamma"'] & & Y\times\mathbb P^1 \arrow[ur,"\alpha"'] &
\end{tikzcd}
\]

Del mismo modo, denotando por $X'$ el subesquema cerrado de $Z'$ obtenido por la explosión de $X=X\times0$ a lo largo de $Y$, es claro que $\mathbb P(N)=X'\cap\widehat N$ se identifica con el hiperplano en el infinito $H$ de $\widehat N$. Las demás notaciones relativas a $X'$ quedan fijadas por el diagrama canónico $(D_2)$.
\[
\begin{tikzcd}[column sep=3.2em,row sep=2.4em,cells={nodes={inner sep=1pt}}]
\text{(D}_2\text{)}\arrow[from=1-3,to=5-2,"g_1"']
& & \widehat N \arrow[rr,"j_1"] & & Z' \arrow[ddd,"f_1"] \\
& H \arrow[ur,"k"] \arrow[rr,"j"'] \arrow[ddd,"g"'] & &
X' \arrow[ur,"v"'] \arrow[ddd,"f"] & \\
& & & & \\
& & & & X\times\mathbb P^1=Z \\
& Y \arrow[rr,"i"'] & & X \arrow[ur,"u"'] &
\end{tikzcd}
\]
